\chapter*{Como ler este livro}
\addcontentsline{toc}{chapter}{Como ler este livro}

O livro não precisa ser lido em ordem. Os capítulos~1--4 são a base comum: os tipos de neurônios, o que calculam os \nt{GLUT} e os \nt{GABA} e em que linguagem eles se comunicam. Depois disso o livro se ramifica. O mapa da fig.~\ref{fig:roadmap} mostra quais capítulos se apoiam em quais: uma seta do capítulo $A$ ao capítulo $B$ significa que $B$ usa conceitos e fórmulas de $A$. O mapa mostra apenas as dependências principais; pequenas referências a capítulos anteriores não atrapalham a leitura.

\begin{figure}[h]
\centering
\begin{tikzpicture}[>=Stealth,scale=0.95,
  ch/.style={draw,thick,rounded corners=3pt,align=center,font=\scriptsize,text width=1.85cm,minimum height=0.62cm,inner sep=2pt},
  base/.style={ch,fill=blue!8}, bus/.style={ch,fill=orange!14}, sum/.style={ch,fill=gray!18},
  io/.style={ch,fill=green!10}, sort/.style={ch,fill=cyan!8}, lab/.style={ch,fill=violet!10},
  dep/.style={->,thick,gray!60!black}]
% foundation
\node[base] (c1) at (0,0) {1 Tipos de neurônios};
\node[base] (c2) at (2.3,0) {2 \nt{GLUT}};
\node[base] (c3) at (4.6,0) {3 \nt{GLUT} e \nt{GABA}};
\node[base] (c4) at (6.9,0) {4 Código Morse};
% broadcast channels and the synthesis
\node[bus] (c5) at (4.6,-1.5) {5 \nt{SERO}};
\node[bus] (c6) at (7.3,-1.5) {6 \nt{DOPA}};
\node[bus] (c7) at (9.2,-0.9) {7 Teorias da dopamina};
\node[bus] (c8) at (9.2,-2.1) {8 \nt{NORA}};
\node[bus] (c9) at (11.5,-2.1) {9 \nt{ACET}};
\node[sum] (c10) at (13.8,-0.9) {10 A máquina inteira};
% inputs and outputs
\node[io] (c12) at (2.3,-4.1) {12 Reconhe-\\cimento};
\node[io] (c11) at (4.6,-4.1) {11 Entradas};
\node[io] (c13) at (6.9,-4.1) {13 Músculos};
\node[io] (c14) at (9.2,-3.05) {14 Dor};
\node[io] (c15) at (9.2,-4.1) {15 Aprender\\movimentos};
\node[io] (c16) at (9.2,-5.15) {16 Química\\do corpo};
% lab, sorting, sleep
\node[lab] (c17) at (11.5,-4.1) {17 Depuração};
\node[lab] (c21) at (13.8,-4.1) {21 Máquinas vivas};
\node[lab] (c22) at (13.8,-5.15) {22 DishBrain};
\node[sort] (c18) at (11.5,-5.15) {18 Neurônios-\\instrumentos};
\node[sort] (c19) at (13.8,-6.2) {19 Número de\\dendritos};
\node[bus] (c20) at (11.5,-6.2) {20 Sono};
% dependencies
\draw[dep] (c1) -- (c2); \draw[dep] (c2) -- (c3); \draw[dep] (c3) -- (c4);
\draw[dep] (c1) -- (c5); \draw[dep] (c4) -- (c6); \draw[dep] (c5) -- (c6);
\draw[dep] (c6) -- (c7); \draw[dep] (c6) -- (c8); \draw[dep] (c8) -- (c9);
\draw[dep] (c7) -- (c10); \draw[dep] (c9) -- (c10);
\draw[dep] (c4) -- (c11); \draw[dep] (c11) -- (c12); \draw[dep] (c11) -- (c13); \draw[dep] (c5) -- (c13);
\draw[dep] (c13) -- (c14); \draw[dep] (c13) -- (c15); \draw[dep] (c13) -- (c16);
\draw[dep] (c15) -- (c17); \draw[dep] (c9) -- (c17); \draw[dep] (c17) -- (c21); \draw[dep] (c21) -- (c22);
\draw[dep] (c16) -- (c18); \draw[dep] (c18) -- (c19); \draw[dep] (c16) -- (c20);
% legend
\begin{scope}[yshift=-7.25cm]
\foreach \s/\t [count=\i] in {base/base, bus/canais e modos, sum/síntese, io/entradas e saídas, sort/classificações de neurônios, lab/laboratório}{
  \pgfmathtruncatemacro{\col}{mod(\i-1,3)}\pgfmathtruncatemacro{\row}{(\i-1)/3}
  \node[\s,text width=0.3cm,minimum height=0.32cm] at ({\col*4.8+0.2},{-\row*0.55}) {};
  \node[anchor=west,font=\scriptsize] at ({\col*4.8+0.55},{-\row*0.55}) {\t};}
\end{scope}
\end{tikzpicture}
\caption{Mapa dos capítulos. A seta vai de um capítulo para aquele que se apoia nele. As demais referências entre capítulos são indicações, não dependências.}
\label{fig:roadmap}
\end{figure}

\section*{Caminhos pelo livro}

\begin{center}
\footnotesize
\setlength{\tabcolsep}{4pt}
\renewcommand{\arraystretch}{1.2}
\begin{tabular}{>{\raggedright}p{3.0cm}>{\raggedright}p{3.9cm}>{\raggedright\arraybackslash}p{6.4cm}}
\toprule
Caminho & Para quem & Capítulos em ordem \\
\midrule
curso de um trimestre & professor, 10 semanas & semana 1: cap.~1--2; 2: cap.~3; 3: cap.~4; 4: cap.~5--6; 5: cap.~7--8; 6: cap.~9; 7: cap.~10; 8: cap.~11--12; 9: cap.~13 e~15; 10: cap.~17 \\
IA e aprendizado de máquina & quem conhece redes neurais e quer entender como o cérebro aprende & 1--4, 5 (leitura rápida), 6, 7, 8, 9, 11 (leitura rápida), 12, 13 (leitura rápida), 15 \\
eletrônica e hardware & projetista de circuitos, desenvolvedor de chips neuromórficos & 1--4, 5, 11, 13, 16, 18, 19, 17 (capítulos 9 e 15 em leitura rápida) \\
neurointerfaces e computadores vivos & quem constrói interfaces cérebro-computador e chips com neurônios vivos & 1, 2, 4, 11, 13, 17 (capítulos 9 e 15 em leitura rápida), 21, 22 (capítulo 12 em leitura rápida) \\
visão geral rápida & engenheiro com um fim de semana livre & 1, 2, 3, 6, 10; as referências do capítulo~6 aos capítulos~4--5 se entendem pelo contexto \\
\bottomrule
\end{tabular}
\end{center}

``Leitura rápida'' significa: basta ler o início do capítulo, as proposições nos quadros amarelos e a tabela final.

No fim de cada capítulo há exercícios: estimativas, deduções, modelagem e projeto; as respostas curtas estão reunidas no apêndice~\ref{sec:answers}. Toda a notação está no apêndice~\ref{sec:notation}, e os experimentos que sustentam as principais afirmações de cada capítulo estão no apêndice~\ref{sec:supplement}. Os números do livro vêm de pequenos programas na pasta \texttt{models/}; eles podem ser executados e modificados.
